\documentclass[11pt]{article}

\usepackage[margin=1.15in]{geometry}
\usepackage{amsmath,amssymb}
\usepackage{booktabs}
\usepackage{graphicx}
\usepackage{listings}
\usepackage{xcolor}
\usepackage[hidelinks]{hyperref}

\lstset{
  basicstyle=\ttfamily\small,
  frame=single,
  framesep=4pt,
  breaklines=true,
  columns=fullflexible,
  xleftmargin=0pt,
  aboveskip=10pt,
  belowskip=10pt,
}

\title{\textbf{Spatio-Temporal Acoustic Monitoring for Predictive
Maintenance: Auditing a Multi-Agent Codebase and Correcting the
Measurements It Reported}}
\author{Matthew Park}
\date{August 2026}

\begin{document}
\maketitle

\begin{abstract}
\noindent
Industrial acoustic monitoring promises to detect mechanical failure before it
occurs, from sound alone. Murmur is a system built for that task: an array of
microphones feeds a spatio-temporal graph neural network over an
inverse-distance acoustic topology, whose per-timestep embeddings drive a
continuous-time liquid network that forecasts time-to-failure, with a
convolutional autoencoder scoring per-node anomalies against each microphone's
own rolling baseline.

This paper reports an audit of that codebase and a substantially corrected
system. The codebase is unusual in one respect that turns out to matter: it was
built by three independent coding agents working concurrently in the same
repository. That produced a specific and reproducible failure pattern. Two
agents implemented acoustic localisation twice and an evaluation harness twice;
one of each survived. Six commits of finished work --- source localisation,
public-dataset benchmarking, anomaly explainability, and a fault taxonomy ---
were stranded on an unmerged branch, and two further modules, webhook alerting
and ONNX export, existed only as untracked files inside an agent worktree and
would have been destroyed by a routine cleanup. An audit performed by one agent
correctly identified eleven defects and, having been placed in a read-only role,
fixed none of them; all eleven were still live.

The most consequential finding is not a crash. The system reported a
time-to-failure $\mathrm{F}_1$ of $1.000$. That figure was an artifact of the
data generator: fault severity was drawn from $U(0.5, 1.0)$ while healthy
sequences sat below $0.1$, leaving a $0.4$-wide interval containing no examples,
with the $0.5$ decision threshold sitting inside it. The classes were separable
by construction. The same draw confined the middle severity band to a sliver,
starving the per-group conformal calibration that exists specifically to hold
coverage in the high-risk strata --- so the one guarantee aimed at the regime an
operator acts in was the one silently disabled.

I close the gap, correct eleven defects with a regression test for each, salvage
and integrate the stranded work, and re-measure. Under the corrected
distribution, $\mathrm{F}_1$ falls from $1.000$ to $0.9672$ and mean absolute
error improves from $0.0608$ to $0.0313$ against a mean-predictor baseline of
$0.2450$. Conformal coverage lands at $0.89$ against a nominal $0.90$ with mean
interval width falling from $0.190$ to $0.119$, and all three severity strata
receive their own radius for the first time, ordered $0.048 < 0.074 < 0.105$
from normal to critical. I do not claim these numbers demonstrate acoustic fault
detection. They are measured on a simulator, and Section~\ref{sec:notclaim}
states precisely what that does and does not establish.
\end{abstract}

\tableofcontents
\newpage

\section{Introduction}

Rotating machinery announces its own failure. A bearing race that has begun to
spall emits a high-frequency squeal with harmonics; a pump drawing below its
required suction head produces broadband noise punctuated by the collapse of
vapour bubbles; a rotor carrying lost balance weight modulates at shaft rate.
These signatures precede functional failure by days or weeks, and a technician
walking the floor with thirty years of experience can often hear them. The
motivation for automating this is not that the task is impossible for humans but
that it does not scale: a plant has thousands of machines and one walkthrough per
shift.

Murmur addresses this with a fixed microphone array and three models in
sequence. Audio is captured in half-second chunks, converted to log-mel
spectrograms on the GPU, and buffered into temporal windows. A spatio-temporal
graph neural network treats the array as a graph --- microphones are nodes,
inverse-distance acoustic coupling gives the edges --- and produces a per-node
embedding at every timestep. Those embeddings drive a closed-form continuous-time
network (CfC) that integrates over real inter-frame intervals to forecast
time-to-failure. In parallel, a convolutional autoencoder scores each frame
against that microphone's own rolling baseline using a median/MAD robust
$z$-score.

\subsection{Motivation for this paper}

This project has an unusual provenance. It was developed by three independent
coding agents operating concurrently on the same repository, in separate git
worktrees, without coordination between them. That arrangement produces a
failure mode worth documenting, because it is not the one intuition predicts.
The naive worry is merge conflicts. What actually happened was quieter:

\begin{itemize}
\item \textbf{Duplicated effort that resolved arbitrarily.} Two agents
independently implemented time-difference-of-arrival source localisation, and
two independently implemented an evaluation harness. In each case one landed on
\texttt{main} and the other did not. The surviving localisation implementation
was not the more sophisticated one; it was the one wired into the inference
worker.

\item \textbf{Work stranded rather than lost.} One agent's branch carried six
commits --- benchmarking against MIMII and DCASE, GCC-PHAT localisation,
anomaly explainability, and a grounded fault taxonomy --- that were never merged.
Two modules, a webhook alerting router and an ONNX export path, were never
committed at all. They existed only as untracked files inside a worktree
directory that a routine \texttt{git worktree remove} would have deleted without
warning.

\item \textbf{An audit that fixed nothing.} A third agent audited the repository
and correctly identified eleven defects. It had been assigned a read-only role
to avoid editing files another session held open. Its report was accurate and
its findings were never applied. All eleven were still live when I verified them
against the tree.
\end{itemize}

None of these are merge conflicts. Version control was working exactly as
designed throughout; the failure is that concurrent agents produce work whose
\emph{integration} nobody owns.

The second motivation is the more important one. The system reported an
$\mathrm{F}_1$ of $1.000$ on time-to-failure classification. A perfect score on
a hard problem should be read as evidence about the measurement rather than the
model, and it was: the label distribution had a hole in it, and the decision
threshold sat in the hole.

\subsection{Contributions}

\begin{enumerate}
\item \textbf{A defect audit} of a concurrently-developed codebase, classifying
eleven verified defects into those that stop the system emitting output, those
that run silently while disabling a documented capability, and those that
corrupt reported measurements (Section~\ref{sec:audit}). Each is characterised
by mechanism and closed with a regression test.

\item \textbf{Correction of the measurement protocol.} The synthetic label
distribution contained a $0.4$-wide interval with no examples, spanning the
classification threshold. I identify the mechanism, close it, and show the
effect on every reported figure (Sections~\ref{sec:data} and \ref{sec:experiments}).

\item \textbf{Repair of the uncertainty quantification.} Per-group conformal
calibration was falling back to a global radius for two of three severity
strata. I show this was a consequence of the same distributional defect
compounded by an undersized calibration split, and restore group-conditional
coverage (Sections~\ref{sec:conformal} and \ref{sec:experiments}).

\item \textbf{Recovery and integration of stranded work.} Roughly 2{,}000 lines
across seven modules were salvaged from an unmerged branch and an uncommitted
worktree, deduplicated against the implementations that had already landed, and
--- critically --- wired into the runtime rather than left as dead code
(Section~\ref{sec:artifact}).

\item \textbf{A verifiable artifact}: 452 automated tests, a Kafka integration
test that genuinely executes rather than skipping, and infrastructure validated
against the tooling that consumes it (Section~\ref{sec:artifact}).
\end{enumerate}

\subsection{What this paper does not claim}
\label{sec:notclaim}

I do not claim that Murmur detects acoustic faults in industrial machinery.
Every number reported here is measured on a synthetic generator, and a
generator is a statement of the author's assumptions rather than evidence about
the world. The corrected generator produces spatially-localised faults with
inverse-square attenuation, per-fault spectral profiles, temporal onset ramps,
and an ambient floor deliberately decoupled from the label --- but a model that
performs well against those assumptions has demonstrated only that it can
recover the structure I put in.

I also do not claim the pre-correction figures were wrong in the sense of being
miscomputed. $\mathrm{F}_1$ was $1.000$; the arithmetic was correct. The claim is
that the quantity being computed did not measure what the surrounding prose said
it measured, and that this is the more dangerous kind of error because nothing
fails.

The corrected and uncorrected numbers are not comparable as model comparisons.
They are measured on different distributions. Where both appear in this paper it
is to quantify the effect of the correction, never to suggest one model is
better than another.

Establishing real detection performance requires the public benchmarks the
repository now contains harnesses for --- MIMII, DCASE, IMS --- run against
recorded machine audio. That is the obvious next step and it is not attempted
here.

\section{Related Work}

\subsection{Acoustic condition monitoring}
Vibration analysis is the established technique for rotating machinery, with
envelope analysis and cepstral methods used to isolate bearing defect
frequencies. Airborne acoustic monitoring trades signal quality for
installation cost: a microphone requires no contact with the machine and can
cover several assets at once, at the price of reverberation, masking, and
ambient noise. The DCASE challenge series established anomalous sound detection
as a benchmark task, and the MIMII corpus provides recordings of pumps, fans,
valves, and slide rails under normal and faulty operation at controlled
signal-to-noise ratios.

\subsection{Unsupervised anomaly detection}
Fault data is scarce by construction --- a well-maintained plant produces
overwhelmingly normal recordings --- so the dominant framing is unsupervised:
model the normal manifold and score departures from it. Reconstruction-error
autoencoders remain the standard baseline. Murmur follows this approach and adds
a per-node robust baseline: each microphone is judged against its own recent
history using a median and median absolute deviation rather than a mean and
standard deviation, because a developing fault contaminates precisely the
statistics used to detect it.

\subsection{Graph neural networks for sensor arrays}
Where sensors have spatial relationships, graph networks exploit them directly.
Spatio-temporal GNNs interleave graph convolution with temporal modelling and
are well established in traffic forecasting. The application here treats
microphones as nodes and acoustic coupling as edges, so that a fault localised
at one machine propagates through the graph with distance-dependent
attenuation.

\subsection{Continuous-time recurrent models}
Liquid time-constant networks and their closed-form approximation (CfC) define
hidden state through an ODE, which makes them natural for irregularly sampled
data: the inter-arrival interval enters the dynamics rather than being assumed
uniform. This matters for an edge deployment where network jitter makes the
sampling interval genuinely variable, and it is the justification for the
architecture choice --- a justification that, as Section~\ref{sec:audit} records,
the original tensor plumbing defeated.

\subsection{Conformal prediction}
Split conformal prediction provides finite-sample marginal coverage without
assumptions on the model or the noise distribution. Mondrian (per-group)
conformal extends this to conditional coverage within strata, which is the
relevant variant here: coverage averaged over a population dominated by healthy
machines can look excellent while being poor exactly where decisions are made.

\section{Defect Audit}
\label{sec:audit}

This section documents defects verified against the tree rather than inferred
from reports. Each was reproduced before being fixed and has a regression test
afterwards; the existing 123-test suite passed throughout, which is the point of
recording them.

\subsection{Class 1: the system stops emitting}

\textbf{A single dead microphone silences the entire array.} The window
assembler required every node to have reported before releasing a graph
snapshot, with no timeout, no eviction, and no degraded mode:

\begin{lstlisting}
def is_complete(self) -> bool:
    if len(self._windows) < self.num_nodes:
        return False
    ...
    return (newest - oldest) <= self.staleness_tolerance
\end{lstlisting}

If one microphone drops out, the remaining three stream indefinitely and nothing
is ever emitted. Worse, the staleness check cannot recover: a node that stops
updating retains its last window forever, so the spread between newest and
oldest never returns below tolerance and even the surviving microphones stop
producing. No metric fired. The dashboard displayed ``waiting for model
predictions'', which is visually identical to a healthy, quiet plant.

For a system whose entire premise is not missing events, going silent is the
worst available failure mode. The assembler now evicts stale windows, releases a
degraded snapshot once a quorum is present after a bounded wait, zero-fills
absent nodes so the topology stays index-aligned, and omits them from telemetry
so nothing is fabricated for a microphone that said nothing.

\subsection{Class 2: silent defects}

These executed without error while disabling a documented capability.

\textbf{A facility-level forecast presented as four machine-level forecasts.}
One time-to-failure value and one embedding were computed from the pooled graph
readout and copied into every node's payload. The dashboard's four per-node
cards were therefore identical by construction --- an operator reading four
independent machine forecasts was reading one number four times. The ST-GNN now
exposes per-node embeddings and each microphone receives a forecast derived from
its own trajectory.

\textbf{A write-only Kafka topic.} Every spectrogram frame from every node was
serialised, compressed, and shipped to a broker that no consumer subscribed to,
with retention paid on all of it. The topic is now opt-in.

\textbf{The rate limiter as a memory exhaustion vector.} Buckets were keyed by
API key or client address --- both attacker-controlled on a public endpoint ---
trimmed within each key but never across keys, and cleared only at shutdown. One
request per spoofed source grew the map without bound. The control added to
prevent denial of service enabled a different one.

\textbf{A tensor cache keyed on memory addresses.} The batched graph topology was
cached under \texttt{edge\_index.data\_ptr()}, which is unique only while the
tensor is alive. Once freed, the allocator may reissue that address to an
unrelated tensor and the cache would return edges belonging to a graph that no
longer exists. The probability is low and the symptom --- a silently wrong
topology --- is severe. The cache now keys on shape and confirms tensor identity,
holding a reference so the address cannot be recycled.

\textbf{An unauthenticated metrics endpoint.} \texttt{/metrics} published
per-node anomaly counts, $z$-scores, and failure forecasts --- a live map of
which machines in the plant are failing --- without authentication even when an
API key was configured, on a LoadBalancer-exposed service.

\subsection{Class 3: measurement defects}

\textbf{The label distribution had a hole in it.} This is the finding that
invalidates a reported number, and it is developed in full in
Section~\ref{sec:data}.

\textbf{At-least-once delivery was overstated.} The module documented that
``offsets advance only after downstream publish succeeds''. The implementation
enqueued messages, served already-completed delivery callbacks, and committed:

\begin{lstlisting}
producer.poll(0)          # serves callbacks that already finished
consumer.commit(...)      # acknowledges data still in the client queue
\end{lstlisting}

Messages still in the librdkafka queue were unacknowledged when the offset
advanced, so a crash in that window loses precisely the frames the comment
promises to protect. Broker-side idempotence guards duplication, not this. The
producer is now drained before offsets move, and a failed drain holds the batch
for redelivery.

\textbf{An integration test that could not fail.} The Kafka fixture probed
availability by constructing a producer and calling \texttt{flush()}. Neither
connects --- librdkafka dials lazily, and flushing an empty queue returns zero
regardless --- so the fixture always reported the broker as present and the test
then failed on a missing message rather than skipping. In continuous
integration the service was moreover started with no configuration at all, and
readiness was probed by opening a TCP socket. An open port is not a broker
serving metadata, so a half-started Kafka would have produced a green run that
exercised nothing.

\textbf{A degenerate $z$-score reaching the operator.} When a rolling baseline
has no spread, the detector multiplied the relative departure by $10^3$,
producing $z$-scores in the tens of thousands that were written to a Prometheus
gauge and interpolated into a language model prompt verbatim. The magnitude is
meaningless --- the ratio is taken against an arbitrarily small median --- so it
now saturates at a documented constant, preserving sign, because a channel that
went quiet is not a bearing failure.

\subsection{Infrastructure}

Three defects would each have prevented deployment and only one was visible.
Bitnami withdrew its versioned tags from Docker Hub, so the Kafka image
referenced by the development compose file, the continuous integration service,
and the production Kubernetes StatefulSet no longer resolves. Only the CI
reference failed loudly; the StatefulSet would have failed at deploy time. All
three now use the Apache project's image, which required translating every
environment variable (the two images use different prefixes) and setting the log
directory explicitly --- the Apache image defaults it to \texttt{/tmp}, so the
StatefulSet would otherwise have retained a healthy 20\,GiB volume claim and
still lost every unconsumed message on restart.

Separately, the test suite could not run in CI at all: a package was added
without being registered for installation, and because CI invokes
\texttt{pytest} directly rather than \texttt{python -m pytest} there is no
working directory on \texttt{sys.path} to fall back on. Six modules failed at
collection, which aborts the entire run --- including the integration job, whose
own test was deselected.

\section{Data and Simulation}
\label{sec:data}

\subsection{The generator}

Each sequence is a window of $50$ frames across $4$ microphones with $64$ mel
bins. An ambient noise floor is drawn per sequence and \emph{independently of
the label}, so a loud healthy machine can out-energise a quiet degrading one;
this removes total loudness as a shortcut. With probability $p$ the sequence
contains a fault, which is assigned a source node, one of three spectral
profiles (bearing, cavitation, imbalance), and a severity. The fault contributes

\[
c_{t,n,b} \;=\; s \cdot r_t \cdot m_t \cdot g_n \cdot \phi_b ,
\qquad
g_n \;=\; \frac{1}{1 + d_{\mathrm{src},n}^{2}} ,
\]

where $s$ is severity, $r_t$ a temporal onset ramp, $m_t$ a slow modulation,
$g_n$ inverse-square attenuation from the source microphone, and $\phi_b$ a
Gaussian envelope over mel bins. The spatial term is what gives the graph
convolution something to learn: a fault at node $0$ is loudest at node $0$ and
attenuates with distance.

\subsection{The defect}

Severity was drawn from $U(0.5, 1.0)$ for faulty sequences, while healthy
sequences received a label from $U(0.0, 0.1)$. Measured over $2{,}000$
sequences, the resulting distribution is:

\begin{center}
\begin{tabular}{lr}
\toprule
Label interval & Count \\
\midrule
$[0.0, 0.1)$ & 1718 \\
$[0.1, 0.5)$ & \textbf{0} \\
$[0.5, 0.6)$ & 64 \\
$[0.6, 0.8)$ & 113 \\
$[0.8, 1.0]$ & 105 \\
\bottomrule
\end{tabular}
\end{center}

The maximum healthy label is $0.0999$ and the minimum faulty label is $0.5011$:
an empty interval $0.4012$ wide. Classification metrics threshold at $0.5$,
which lies inside it. The two classes were linearly separable with a margin, and
$\mathrm{F}_1 = 1.000$ followed for any model that could determine whether a
fault signature was present at all --- a far easier question than estimating how
advanced it is.

This also distorted the severity strata used for conformal grouping. With bands
at $[0, 0.33)$, $[0.33, 0.66)$, $[0.66, 1]$, the faulty draw could only reach
\emph{warning} through the sliver $[0.5, 0.66)$, giving $5.0\%$ of sequences
against $85.9\%$ normal.

\subsection{The correction}

Severity is now drawn from $U(0.05, 1.0)$. This places examples on both sides of
the decision threshold and overlaps the healthy band at the bottom, so a faintly
degrading machine is genuinely difficult to distinguish from a healthy one ---
which is the regime predictive maintenance operates in. The fault rate rises to
$0.40$ and $N$ to $4{,}000$; both are properties of the simulator chosen so that
per-stratum calibration is estimable, not claims about real fleets.

\section{Method}

The pipeline is four stages. Log-mel spectrograms are computed on GPU in
micro-batches and buffered into per-node windows of $50$ frames.

\textbf{Spatio-temporal GNN.} Per-microphone temporal self-attention with
sinusoidal positional encoding runs first, so each microphone attends over its
own history without leaking across the array. Attention is permutation-invariant
without positional encoding, which for a model detecting temporal signatures is
a correctness defect rather than a refinement. Graph convolution then runs at
\emph{every} timestep --- $B \times S$ disjoint copies of the topology in one
call --- which is what makes the spatial convolution time-resolved, and is the
dominant training cost. Readout produces both a pooled graph embedding and
per-node embeddings ($0.31$\,M parameters).

\textbf{Liquid network.} A CfC integrates the per-node embedding sequence over
the measured inter-frame intervals ($0.05$\,M parameters). The sequence form is
essential and was the subject of an earlier defect: feeding a single pooled
vector repeated across time gives a continuous-time model a constant signal and
discards the entire justification for the architecture.

\textbf{Anomaly scoring.} A convolutional autoencoder reconstructs single frames;
per-node reconstruction error is converted to a robust $z$-score against that
node's rolling median and MAD. Group normalisation rather than batch
normalisation, because the worker scores one frame at a time.

\textbf{Grounded diagnosis.} For flagged frames the per-bin reconstruction error
is decomposed across frequency bands. Because the anomaly score \emph{is} the
mean squared reconstruction error, this decomposition is exact rather than an
approximation of sensitivity: each band's share is literally its contribution.
Bands are matched against a fault catalogue using a Bhattacharyya coefficient
over normalised distributions, which prevents a broad signature (cavitation,
$1$--$8$\,kHz) from winning every high-frequency diagnosis by covering more
ground than a narrow one.

\section{Uncertainty Quantification}
\label{sec:conformal}

The forecaster emits a sigmoid. Nothing in the training objective makes $0.7$
mean ``fails 70\% of the time'', and that is the number a maintenance planner
would schedule an outage against. Split conformal prediction converts it to an
interval with finite-sample marginal coverage without assumptions on the model
or error distribution: hold out a calibration set, measure realised residuals,
take a quantile.

Grouping matters more than the base method here. Marginal coverage averaged over
a population that is $70\%$ healthy can be excellent while coverage among
critical forecasts is poor --- and critical forecasts are the only ones anyone
acts on. Mondrian conformal calibrates within strata, requiring
\texttt{min\_group\_size} samples per stratum before a group receives its own
radius.

Two compounding defects prevented this from working. The label gap starved the
middle band, and the calibration set --- half of a $15\%$ test split --- was
$75$ samples. The \emph{warning} stratum received $8$. Both high-risk bands fell
back to the global radius, so the guarantee held on the healthy majority and
nowhere else. Enlarging the test split to $20\%$ raises calibration to $400$
samples; combined with the corrected label distribution, all three strata now
qualify.

\section{Experiments}
\label{sec:experiments}

\subsection{Protocol}

Identical architecture, optimiser, schedule, and seed ($1337$) throughout; $50$
epochs with cosine annealing and early stopping on validation loss. The
calibration set is half the test split, disjoint from both training and
validation --- the validation set is not exchangeable with unseen data because
early stopping selected on it. Coverage is measured on the other half, which
neither the model nor the calibrator has seen.

\subsection{Effect of the correction}

Run A is the original configuration. Run B closes the label gap and raises $N$.
Run C additionally enlarges the calibration split.

\begin{center}
\begin{tabular}{lccc}
\toprule
& \textbf{A} (original) & \textbf{B} (gap closed) & \textbf{C} (final) \\
\midrule
Sequences $N$              & 1000   & 4000   & 4000   \\
Train/val/test             & 70/15/15 & 70/15/15 & 65/15/20 \\
\midrule
$\mathrm{F}_1$             & 1.0000 & 0.9712 & 0.9672 \\
Precision                  & 1.0000 & 0.9593 & 0.9529 \\
Recall                     & 1.0000 & 0.9833 & 0.9818 \\
MAE                        & 0.0608 & 0.0375 & 0.0313 \\
Baseline MAE               & 0.3255 & 0.2433 & 0.2450 \\
\midrule
Coverage (nominal $0.90$)  & 0.933  & 0.890  & 0.890  \\
Mean interval width        & 0.190  & 0.141  & 0.119  \\
Calibration $n$            & 75     & 300    & 400    \\
Strata with own radius     & 1 / 3  & 2 / 3  & \textbf{3 / 3} \\
\bottomrule
\end{tabular}
\end{center}

\subsubsection{Interpretation}

$\mathrm{F}_1$ falling from $1.000$ to $0.9672$ is the headline result and it is
an improvement in the measurement, not a regression in the model. Run A's
perfect score reported that no example lay near the threshold. Run C's precision
of $0.9529$ and recall of $0.9818$ describe a detector that now makes both false
positives and misses on a task where the classes overlap --- and that is a
quantity with content.

MAE improves in the same direction the classification metric worsens
($0.0608 \rightarrow 0.0313$), which is initially counterintuitive and is
explained by the baseline. The mean-predictor baseline also falls
($0.3255 \rightarrow 0.2450$) because the label distribution is less extreme;
the ratio of model to baseline improves from $5.4\times$ to $7.8\times$. The
regression task became easier while the thresholded classification task became
harder, because filling the gap moves mass toward the middle in both cases.

Coverage moves from $0.933$ to $0.890$ against a nominal $0.900$. Run A was
\emph{over}-covering by $3.3$ points, which is not a success: it means intervals
were wider than the data required, and width is the cost an operator pays.
Run C is within one point of nominal with intervals $37\%$ narrower.

\subsection{Group-conditional calibration}

The fitted radii in Run C:

\begin{center}
\begin{tabular}{lrr}
\toprule
Stratum & Calibration $n$ & Radius \\
\midrule
normal   & 293 & 0.0478 \\
warning  &  44 & 0.0741 \\
critical &  63 & 0.1054 \\
\midrule
global   & 400 & 0.0513 \\
\bottomrule
\end{tabular}
\end{center}

The ordering is the result worth noting. Uncertainty grows monotonically with
severity: the model is precise about healthy machines and progressively less so
as failure approaches. That is a physically sensible structure, and it is
exactly what the grouping exists to discover --- and could not previously
express, because in Runs A and B the high-risk strata inherited a global radius
of $0.051$, roughly half the $0.105$ that critical forecasts actually require.
Under the original configuration, intervals on the most consequential
predictions were approximately twice too narrow, and the marginal coverage
figure of $0.933$ concealed it.

\subsection{An intermediate hypothesis that was wrong}

Run B left the \emph{warning} stratum at $26$ samples, still below the threshold
of $30$. My initial explanation was that a sigmoid output trained under MSE
compresses predictions toward the extremes, under-populating the middle band in
prediction space --- groups are formed on predictions, not labels, since labels
are unavailable at inference. Measuring the bucket occupancy of the two
distributions refuted this: labels gave $440/70/90$ across normal, warning, and
critical, and predictions gave $440/69/91$. No compression exists. The shortfall
was ordinary sampling variance in a $300$-sample calibration half, and the
correct remedy was a larger calibration set rather than a different model. I
record this because the plausible mechanism and the actual one prescribed
different fixes.

\section{The Artifact}
\label{sec:artifact}

\subsection{Recovered work}

Roughly $2{,}000$ lines were salvaged from an unmerged branch and an
uncommitted worktree: a benchmark harness for MIMII, DCASE, and IMS; performance
benchmarks; a scenario generator with lead-time events; spectral anomaly
explainability; a fault taxonomy; a webhook alerting router with cooldown
suppression; and ONNX export with INT8 quantisation.

Duplicates were resolved rather than merged. The stranded localisation module
was discarded in favour of the implementation already wired into the worker; two
independent ROC implementations were collapsed to one while preserving the
benchmark variant's deliberate policy of reporting chance rather than raising on
a degenerate split.

Salvage alone would have reproduced the defect the audit identified --- modules
that exist but nothing calls. The recovered components are integrated: the
worker computes a spectral explanation for flagged frames, matches it against
the taxonomy, and attaches both to telemetry, where the API states them in the
generation prompt and in the templated fallback. This directly addresses a
finding that the projector conditioning generated text on audio is a random
projection until trained, so unconditioned prose was being presented as
diagnosis. Alerting routes to Slack, PagerDuty, or a webhook, with all sinks
empty by default: a monitoring system should not page anyone until asked.

One defect was found in the recovered code itself. The ONNX verification
routine compared the exported graph against a \emph{freshly initialised} model,
so it could only pass when the tolerance was wide enough to be meaningless.

\subsection{Verification}

452 tests, executing in roughly 25 seconds. Every audit finding has a regression
test, because all eleven were invisible to the suite that existed. Kafka
integration is verified against a real broker rather than skipped: the readiness
probe performs the same metadata round trip the test fixture does, so passing it
means the test genuinely runs.

Infrastructure is validated against the tooling that consumes it rather than by
inspection --- manifests through \texttt{kubeconform} ($14$ resources, $0$
invalid), the container image pulled and its user, healthcheck path, and default
log directory confirmed directly, and the package layout checked by reproducing
the Dockerfile's copy set in isolation.

\section{Discussion}

\subsection{On perfect scores}

The practical lesson is narrow and transferable: a perfect score is a
measurement result, and the first hypothesis should concern the measurement. The
defect here required no unusual insight to find --- a histogram of the labels
shows it immediately --- but nothing prompted anyone to draw one, because the
number was good and good numbers do not invite scrutiny. The system had a
comprehensive test suite, and it passed. No test asserted anything about the
distribution the tests were run against.

\subsection{On concurrent agents}

The failure pattern from three agents in one repository was not conflicting
edits. It was unowned integration: duplicated work resolved by which branch
happened to merge, finished work stranded on branches nobody merged, valuable
work existing only as untracked files, and an accurate audit whose findings
nobody applied. Each agent behaved reasonably in isolation. Every one of these
is a coordination failure at the boundary rather than a defect inside any
agent's work, and version control did not detect any of them because none is a
conflict.

\subsection{Limitations}

Results are measured on a simulator and establish only that the pipeline
recovers structure that was deliberately placed in it. The fault catalogue is
small and its frequency bands are nominal rather than derived from
machine-specific kinematics. The array is four microphones in one room; the
graph is small enough that the spatial convolution's contribution over a
per-node MLP has not been isolated by ablation. The conformal guarantee assumes
exchangeability between calibration and production, which acoustic drift breaks
--- a drift monitor exists to catch coverage collapse, but the guarantee is
conditional on it. Alert escalation deliberately bypasses the suppression
cooldown, so a signal oscillating between warning and critical will re-page on
each upward transition; this is defensible but untested against real fault
progressions.

\section{Conclusion}

I audited a spatio-temporal acoustic monitoring system built concurrently by
three independent agents, verified eleven defects that a previous audit had
correctly identified and left unfixed, and corrected them with regression tests.
The most consequential finding was not a crash but a measurement: a reported
$\mathrm{F}_1$ of $1.000$ that reflected a $0.4$-wide hole in the label
distribution spanning the decision threshold. Closing it moved $\mathrm{F}_1$ to
$0.9672$ and, together with an enlarged calibration split, restored
group-conditional conformal coverage across all three severity strata for the
first time --- revealing that intervals on critical forecasts had been
approximately twice too narrow, concealed behind an apparently healthy marginal
coverage figure.

The system does not yet demonstrate acoustic fault detection, and this paper
does not claim it does. What it provides is a pipeline in which such a claim
could be honestly evaluated: the measurement protocol is corrected, the
uncertainty quantification does what it advertises in the regime that matters,
and the benchmark harnesses against recorded machine audio are present and
tested. Running them is the next step.

\end{document}
